\documentclass[11pt,a4paper]{article}
\usepackage[margin=21mm,headheight=15pt]{geometry}
\usepackage[T1]{fontenc}
\usepackage{lmodern,amsmath,amssymb,booktabs,tabularx,enumitem,xcolor,fancyhdr}
\usepackage[hidelinks]{hyperref}
\definecolor{accent}{RGB}{0,114,178}
\setlength{\parindent}{0pt}
\setlength{\parskip}{6pt}
\setlist{nosep,leftmargin=*}
\pagestyle{fancy}
\fancyhf{}
\fancyhead[L]{\small Econometrics 25117 | Instructor guide}
\fancyhead[R]{\small 2026--2027}
\fancyfoot[L]{\small Private teaching notes}
\fancyfoot[R]{\thepage}
\newcommand{\heading}[1]{\section*{\color{accent}#1}}
\newcommand{\slideheading}[1]{\subsection*{\color{accent}#1}}
\newcommand{\cue}[2]{\par\noindent\textbf{#1}\quad #2\par}
\newcommand{\slideguide}[3]{\slideheading{Slide #1 | #2}\cue{Guide}{#3}}
\setlength{\emergencystretch}{1em}
\newcommand{\E}{\mathbb{E}}
\newcommand{\Var}{\operatorname{Var}}
\newcommand{\Cov}{\operatorname{Cov}}
\begin{document}
\begin{center}{\Large\bfseries Multiple regression and omitted-variable bias}\end{center}
\textbf{Slide route:} Lecture5v2.tex: OVB, multiple regression, goodness of fit, multicollinearity, controls and conditional mean independence. Chapter 6.
\heading{Session 7 | Comparing like with like}
\textbf{Outcomes:} Predict the sign of omitted-variable bias, interpret a partial coefficient, distinguish perfect from high multicollinearity, and choose a reference category.
\heading{100-minute route}
\begin{tabularx}{\textwidth}{@{}rX@{}}\toprule
0--10 & Post-midterm conceptual retrieval.\\
10--30 & OVB slides and education/ability board example.\\
30--45 & Multiple regression and holding controls fixed.\\
45--55 & Goodness of fit and pause.\\
55--75 & Multicollinearity and the dummy-variable trap.\\
75--90 & Control variables and conditional mean independence.\\
90--100 & Exit questions and next topic.\\\bottomrule
\end{tabularx}
\heading{Slide-by-slide script}
\slideguide{1}{Title}{Open by asking what changes when a regression compares units that differ in several ways.}
\slideguide{2}{Roadmap}{Preview OVB, multiple regression, fit, multicollinearity, and controls. Keep the phrase ``holding fixed'' visible.}
\slideguide{3}{Last lesson}{Retrieve the simple-regression slope and ask what makes its causal interpretation fail.}
\slideguide{4}{Last lesson}{Use the class-size example to distinguish a partial relationship from an unconditional comparison.}
\slideguide{5}{OVB}{Introduce an omitted ability variable. Ask students to predict the sign before seeing the formula.}
\slideguide{6}{OVB}{Write the two links required for bias: ability affects wages and ability is correlated with education.}
\slideguide{7}{OVB}{Work $\text{simple slope}=.06+.40(.10)=.10$. Circle the bias component $.04$.}
\slideguide{8}{OVB}{Explain that the sign rule depends on both links. If either link is zero, this simple omitted-variable bias disappears.}
\slideguide{9}{OVB}{Ask whether adding a control automatically creates a causal estimate. No; the control can be measured poorly or be endogenous.}
\slideguide{10}{OVB}{Use a diagram with education, ability, and wage. Have students point to the backdoor path.}
\slideguide{11}{Identifying causal effects}{Return to the intervention question: what would happen to one person under a different education level?}
\slideguide{12}{Identifying causal effects}{Contrast observed comparisons with comparisons holding relevant determinants fixed.}
\slideguide{13}{Another way to see this}{Use residualization language: remove the part of education explained by controls, then compare the remaining education variation.}
\slideguide{14}{Another way to see this}{Warn that partialling out is an algebraic equivalence, not a guarantee that the controls are valid.}
\slideguide{15}{What to do about OVB}{Ask students to propose randomization, a credible instrument, or a defensible control strategy.}
\slideguide{16}{Multiple regression model}{Write $Y=\beta_0+\beta_1X_1+\beta_2X_2+u$ and say exactly what $\beta_1$ holds fixed.}
\slideguide{17}{Multiple regression model}{Use a two-regressor prediction example and identify which coefficient changes when the second regressor is added.}
\slideguide{18}{Example for $k=2$}{Calculate one fitted value from the displayed coefficients. Ask students to keep the intercept and both regressors.}
\slideguide{19}{Example for $k=2$}{Interpret a coefficient as a one-unit change in one regressor at fixed values of the others.}
\slideguide{20}{OLS estimator in multiple regression}{Connect the normal equations to orthogonality with every included regressor.}
\slideguide{21}{OLS estimator in multiple regression}{Ask whether residual orthogonality proves exogeneity. No; it is imposed by the minimization.}
\slideguide{22}{Subsidized meals example}{Read the outcome, treatment, controls, and sample before interpreting the coefficient.}
\slideguide{23}{Subsidized meals example}{Compare a short and long specification. Ask which coefficient changes and what that suggests about confounding.}
\slideguide{24}{Goodness of fit}{Show that adding regressors cannot lower in-sample $R^2$ on the same observations.}
\slideguide{25}{Goodness of fit}{Warn that a higher $R^2$ is not evidence of a better causal model.}
\slideguide{26}{Goodness of fit}{Use adjusted $R^2$ only as a fit summary, not as a substitute for theory or design.}
\slideguide{27}{OLS assumptions for causal inference}{List random sampling, variation, conditional mean zero, and finite moments. Ask which assumption is causal.}
\slideguide{28}{Multicollinearity}{Distinguish perfect collinearity, which prevents separate estimation, from high collinearity, which inflates uncertainty.}
\slideguide{29}{Perfect multicollinearity}{Use $X_3=X_1+X_2$ as a quick algebraic example. Ask which coefficient is not separately identified.}
\slideguide{30}{High multicollinearity}{Explain that large standard errors do not imply bias. Precision and identification are different.}
\slideguide{31}{Dummy-variable trap}{For categories A, B, and C with an intercept, drop one indicator. Ask students to name the reference group.}
\slideguide{32}{Omitted category}{Use the fitted means 50, 55, and 47 from this guide. Change the reference category and show that fitted means stay the same.}
\slideguide{33}{Control variables}{Ask whether a proposed control is pre-treatment, a confounder, a mediator, or a collider.}
\slideguide{34}{Conditional mean independence}{Interpret $\E(u\mid X_1,X_2)=0$ as no systematic remaining error after conditioning on both regressors.}
\slideguide{35}{Causal assumptions}{Return to the wage example and name the assumption that ability is now adequately handled.}
\slideguide{36}{Summary}{Have students explain OVB, partial coefficients, and multicollinearity in one sentence each.}
\slideguide{37}{Material}{Assign one specification-comparison exercise and ask students to prepare a joint null for the next class.}

\heading{Worked example | The sign and size of OVB}
\cue{Use with slides}{Lecture5v2 slides 5--10: Omitted Variable Bias. Put the education/ability calculation on the board before slide 6 if you want them to predict the sign, then reveal the size calculation as the slide logic develops.}
Let the synthetic population model be
\[
\log(wage)=\beta_0+.06\,education+.40\,ability+u,
\]
with $\E(u\mid education,ability)=0$. Suppose
$\Cov(education,ability)/\Var(education)=.10$. Regressing log wage on education alone targets
\[
.06+.40(.10)=.10.
\]
The omitted ability term makes the simple-regression slope exceed the partial effect by $.04$.
\cue{Ask before the arithmetic}{Which two conditions create OVB? The omitted variable must affect the outcome conditional on included regressors, and it must be correlated with the included regressor after accounting for the model's other included controls.}
\cue{Interpretation}{Holding ability fixed, the model implies an approximately 6\% wage change per additional education year, with an exact model-implied multiplicative change of $e^{.06}-1\simeq6.18\%$. Log interpretation will be developed in session 9.}

\newpage
\heading{Side examples | Controls and collinearity}
\cue{Use with slides}{Lecture5v2 slides 16--19 for partialling-out intuition, slides 26--32 for multicollinearity and the dummy-variable trap, and slides 33--35 for control variables and conditional mean independence.}
\cue{Partialling-out intuition}{First remove the part of education predicted by the controls. Also remove the part of the outcome predicted by the controls. Regress the outcome residual on the education residual. The slope is the education coefficient from the full OLS regression using the same observations. This is a teaching intuition, not an additional assumption.}
\cue{Dummy trap}{For three exhaustive categories A, B, C, an intercept and all three indicators satisfy $1=D_A+D_B+D_C$. Drop one indicator or the intercept. With A as the reference, an intercept of 50 and coefficients 5 on B and $-3$ on C imply fitted means 50, 55, and 47.}
\cue{High versus perfect collinearity}{Perfect collinearity prevents unique identification of all coefficients. High but imperfect collinearity can inflate SEs; it does not by itself bias OLS under the other assumptions.}
\cue{Goodness of fit}{Adding a regressor cannot lower in-sample $R^2$ on the same sample. Adjusted $R^2$ can fall. Neither chooses valid causal controls automatically.}
\heading{Teaching comments}
Do not say that including every available variable always improves causal inference. A post-treatment mediator can remove part of the effect of interest; a collider can introduce bias. At this level, establish the question and timing of variables before deciding what to control for.

For a treatment effect, conditional mean independence can justify the coefficient on treatment with suitable controls without making every control coefficient a causal effect. Emphasize this distinction in the final control-variable slides.
\heading{Exit ticket and answers}
\cue{1}{Omitted $Z$ increases $Y$ and is negatively correlated with $X$. Bias in the simple $X$ coefficient? Negative.}
\cue{2}{Does a large SE caused by collinearity prove no relationship? No; precision is weak.}
\cue{3}{If the reference category changes, do fitted category means change? No; the parameterization changes.}
\cue{If short / next preparation}{Shorten repeated subsidized-meal tables, keeping one comparison of specifications. Read Chapter 7 and prepare to test two coefficients together.}
\textbf{Source:} Lecture5v2.tex. Numerical examples are synthetic.

\end{document}
